\chapter{Γλωσσάρι όρων}
\label{app:glossary}

\begin{table}[H]
  \centering
  \begin{tabular}{ll}
    \hline
    Ελληνικός όρος & Αγγλικός όρος \\
    \hline
    διτονική ακολουθία & \tl{bitonic sequence} \\
    διτονική ταξινόμηση & \tl{bitonic sort} \\
    περικομμένη διτονική ταξινόμηση & \tl{truncated bitonic sort} \\
    χρονοδιάγραμμα επιπέδων & \tl{layer schedule} \\
    συγχώνευση επιπέδων & \tl{layer fusion} \\
    δίκτυο ταξινόμησης & \tl{sorting network} \\
    σύγκριση-ανταλλαγή & \tl{compare-exchange} \\
    σωρός & \tl{heap} \\
    επιλογή με βάση την αναπαράσταση & \tl{radix selection} \\
    ψηφίδα & \tl{die} \\
    πλακίδιο & \tl{tile} \\
    πλακίδιο διασύνδεσης & \tl{shim tile} \\
    κοινή μνήμη & \tl{shared memory} \\
    απομονωτής & \tl{buffer} \\
    διπλή ενταμίευση & \tl{double buffering} \\
    άμεση προσπέλαση μνήμης & \tl{direct memory access (DMA)} \\
    νήμα & \tl{thread} \\
    διανυσματοποίηση & \tl{vectorization} \\
    βαθμωτός & \tl{scalar} \\
    ροή δεδομένων & \tl{dataflow} \\
    διοχέτευση & \tl{pipelining} \\
    υποδομή μετρήσεων & \tl{measurement harness} \\
    ταβάνι & \tl{(roofline) ceiling} \\
    υπολογιστική ένταση & \tl{operational intensity} \\
    υλοποίηση αναφοράς & \tl{reference implementation} \\
    υλοποίηση επαλήθευσης & \tl{ground truth} \\
    κίνηση δεδομένων & \tl{data traffic} \\
    εύρος ζώνης μνήμης & \tl{memory bandwidth} \\
    γινόμενο ενέργειας-καθυστέρησης & \tl{energy-delay product} \\
    κατανομή εισόδου & \tl{input distribution} \\
    \hline
  \end{tabular}
  \caption{Αντιστοίχιση ελληνικών και αγγλικών όρων.}
  \label{tab:glossary}
\end{table}
